\begin{frame}{Multi-query Diagnosis of Deep Self-localization Network via Reciprocal Rank Fusion}
\scriptsize
\vspace{-0.8em}

\begin{table}[h!]
\centering
\renewcommand{\arraystretch}{1.5} % increases row height slightly more
\setlength{\tabcolsep}{2.5pt}     % reduces horizontal padding a bit more

\begin{tabular}{|p{2.6cm}|p{4.3cm}|p{4.3cm}|}
\hline
\textbf{Citation} & \textbf{Summary} & \textbf{Advantages \& Disadvantages} \\
\hline
\tiny\textit{[7] M. Yuudai, H. Tomoe, and T. Kanji, 2022 IEEE/SICE Int. Symp. Syst. Integr. (SII), Narvik, pp. 627–632, 2022.
DOI: 10.1109/SII52469.2022.970 \newline 8609.}
&
Proposes a \textbf{DSL} method to diagnose performance drops in \textbf{CNN}-based visual place classifiers under changing environmental conditions. Uses multiple diagnostic queries with image masking and combines results via \textbf{RRF} to identify regions affected by domain shifts, enabling targeted retraining instead of full system retraining.
&
\textbf{Advantages:}
\begin{itemize}\itemsep0.25em
\item Resource-efficient; retrains only affected regions.
\item Practical; uses standard visual data without extra sensors.
\item Multi-query fusion improves diagnostic reliability.
\end{itemize}

\textbf{Disadvantages:}
\begin{itemize}\itemsep0.25em
\item Limited masking techniques.
\item Simplistic place partitioning; no advanced clustering.
\item Tested only on VGG16; generalizability unclear.
\item Additional computational cost from multiple queries.
\end{itemize} \\
\hline
\end{tabular}
\end{table}
\end{frame}
